\title{\textbf{Higher Quadratic Spectral Jets}\\[7pt]
\large Harmonic--Stieltjes factorization defects,\\
polylogarithmic second derivatives, and exact mass collisions}
\author{A research continuation prepared for Vladimir Reshetnikov's ProveIt programme\\
\small Mathematical derivations and computational audit with OpenAI assistance}
\date{11 October 2026}
\begin{document}
\maketitle
\thispagestyle{plain}
\begin{abstract}
The first derivative of a zeta-regularized quadratic spectrum can admit an
anomaly-free factorization even when none of its higher derivatives does.
We determine the exact distinction. For every polynomial weight and every
spectral derivative order, an explicit harmonic-number coefficient formula
separates normally convergent sums from Hurwitz derivatives and generalized
Stieltjes constants. A finite local term records the collision with the
Hurwitz pole. Even weights annihilate all such local residues, but no
nonzero polynomial weight annihilates a fixed higher-order defect for every
positive mass. A joint meromorphic formula explains the discontinuous
first-derivative behavior at noninteger-weight resonances.

For the bilateral quadratic spectrum we evaluate the second derivative at
every nonpositive integer using finitely many polylogarithms, their order
derivatives, and one explicitly defined convergent polylogarithm integral
family. A Gamma normalization then gives an exact mass-antiderivative
ladder at every spectral order, with complete harmonic coefficients and
fixed boundary constants. For two distinct quadratic masses we give the
second derivative for an arbitrary polynomial weight, its fully explicit
polynomial subtraction, an ordinary one-dimensional polylogarithmic
integral, and a convergent collision expansion whose every coefficient is
a finite polygamma expression. A separate logarithmic-square identity
provides Gamma-weighted integrals and a dilogarithmic antiderivative.

These results address four concrete extension questions in the preceding
\emph{Quadratic-Spectrum Stieltjes Identities} report. The classical
multiplicative-anomaly and Bessel ingredients retain their attribution.
No minimality of the added integral coordinate, numerical period
independence, or resolution of the canonical Gaussian conjectures is
asserted. The delivery includes all-index proofs, 834 finite exact
assertions, five rejected mutation controls, 63 independent numerical
comparisons, source provenance, and additive integration notes.
\end{abstract}
\medskip
\noindent\textbf{Source anchor.} Repository metadata were read at \pin.
The matching-title predecessor TeX was read from the user's Library;
its identity with the incoming ZIP bytes was not established. No remote
repository files were changed.
\clearpage
\begingroup\small\tableofcontents\endgroup
\clearpage
